\chapter{Operating Procedure and Radio Terminology}
\label{ch:operating-procedures}

Good operating practice is what makes a shared, crowded radio spectrum
usable by everyone. This chapter defines the core vocabulary and
procedural concepts used universally across amateur radio operating,
independent of any particular country's specific rules.

\section{Core Definitions}

\begin{itemize}
\item \textbf{Station}: one or more transmitters, receivers, or
transceivers, together with any associated equipment, necessary at one
location to carry out a radiocommunication service.
\item \textbf{Emission}: the radiation produced by a transmitting
station --- the actual radio energy sent out, as distinct from the
information it may carry.
\item \textbf{Radio spectrum}: the range of electromagnetic frequencies
used for radiocommunication, allocated internationally by the ITU
(Chapter~\ref{ch:regulatory-bodies}) and subdivided nationally among
different services.
\item \textbf{Interference}: the effect of unwanted energy, from one or
more emissions or from other sources, on the reception of a wanted
signal --- ranging from mild degradation to complete disruption.
\end{itemize}

\section{Emission Classification (Emission Designators)}

Internationally, radio emissions are described using a standardized
alphanumeric code (an \textbf{emission designator}) capturing bandwidth,
modulation type, nature of the modulating signal, and type of
information transmitted. The general structure, defined by the ITU
Radio Regulations, is:
\[
\underbrace{\text{Bandwidth}}_{\text{e.g., 2K80}}\ \underbrace{\text{Letter}}_{\text{modulation type}}\ \underbrace{\text{Digit}}_{\text{modulating signal}}\ \underbrace{\text{Letter}}_{\text{information type}}
\]

\begin{center}
\begin{tabular}{@{}ll@{}}
\toprule
Symbol (1st letter, modulation) & Meaning \\
\midrule
N & Unmodulated carrier \\
A & Double-sideband AM \\
J & Single-sideband, suppressed carrier (SSB) \\
F & Frequency modulation \\
G & Phase modulation \\
\bottomrule
\end{tabular}
\quad
\begin{tabular}{@{}ll@{}}
\toprule
Symbol (3rd, information) & Meaning \\
\midrule
A & Telegraphy, aural (Morse by ear) \\
B & Telegraphy, automatic reception \\
E & Telephony (voice) \\
F & Television/video \\
W & Combination of the above \\
\bottomrule
\end{tabular}
\end{center}

\begin{examplebox}[Reading an emission designator]
Interpret the emission designator \texttt{J3E}.

\emph{Solution.} J = single-sideband, suppressed carrier; 3 = a single
analog channel containing quantized or digital information... (for
voice, simply a single analog channel); E = telephony. \texttt{J3E} is
the standard designator for an SSB voice signal --- exactly what most
amateur SSB voice operation uses.
\end{examplebox}

\begin{examtip}
You do not need to memorize every possible emission designator
combination, but recognizing the common ones --- \texttt{A1A} (CW,
on-off keyed Morse), \texttt{J3E} (SSB voice), \texttt{F3E} (FM voice),
\texttt{F1B}/\texttt{J2B} (common digital mode designators) --- appears
frequently on exams worldwide, since this designator system is one of
the genuinely universal, non-country-specific pieces of amateur radio
terminology.
\end{examtip}

\section{Simplex and Duplex Operation}

\begin{itemize}
\item \textbf{Simplex}: both stations transmit and receive on the same
frequency, taking turns (one station transmits while the other
listens, then they swap) --- the default mode for most HF and casual
VHF/UHF contacts.
\item \textbf{Half-duplex}: communication can occur in both directions,
but not simultaneously; a push-to-talk (PTT) transceiver used with
separate transmit and receive frequencies (as on many VHF/UHF
repeaters) is technically half-duplex from the human operator's
perspective, even though the underlying repeater itself transmits and
receives simultaneously on different frequencies.
\item \textbf{Full duplex}: both directions of communication happen
simultaneously, each on its own frequency, with both parties able to
transmit and receive at the same time (as in a telephone call) --- less
common in amateur voice operation but used in some satellite and
repeater-linking contexts, and by the repeater itself even when users
experience it as half-duplex.
\end{itemize}

\section{Repeater Operation}

A \textbf{repeater} receives a signal on one frequency (the
\textbf{input}) and simultaneously retransmits it, usually with higher
power and from an elevated, well-sited antenna, on a different
frequency (the \textbf{output}), extending the effective range of
low-power mobile and handheld stations well beyond normal line-of-sight
limits. The frequency difference between input and output is the
\textbf{offset} (a fixed, band-dependent value in most regions,
occasionally requiring a specific value that differs by country or
region). Many repeaters use a sub-audible tone or digital code
(\textbf{CTCSS} or \textbf{DCS}) transmitted alongside the voice signal
to selectively open the repeater only for users sending the correct
tone/code, reducing accidental keying from unrelated signals on the
input frequency.

\section{Standard Operating Conventions}

\begin{itemize}
\item \textbf{Calling frequency / calling channel}: a widely recognized
frequency used to establish initial contact before moving to a working
frequency, common in VHF/UHF simplex and in some HF sub-bands.
\item \textbf{CQ call}: a general call inviting any station to respond
(``CQ, CQ, CQ, this is \dots''), as opposed to calling a specific
station directly.
\item \textbf{Break-in / QSK}: full break-in CW operation, where the
transmitter momentarily mutes between individual Morse elements so the
operator can hear (and, if needed, immediately interrupt) the other
station even while their own signal is going out.
\item \textbf{Roger, over, out}: procedural words indicating,
respectively, that a transmission was received and understood, that the
operator is finished transmitting and awaiting a reply, and that the
contact (or that operator's part of it) is complete and no reply is
expected.
\end{itemize}

\begin{examtip}
Standard procedural words and Q-codes (Chapter~\ref{ch:phonetics-codes})
exist precisely because they are unambiguous across languages and
accents, unlike improvised phrasing --- using them correctly is both
good practice and frequently tested on exams.
\end{examtip}

\section{Interference and Good Spectrum Citizenship}

Because spectrum is shared, every operator has a responsibility to
minimize unnecessary interference: listening before transmitting
(especially before calling CQ or joining a frequency already in use),
using no more power than needed for reliable communication, keeping
transmitted bandwidth as narrow as the mode reasonably requires
(Chapter~\ref{ch:non-sine}), and promptly addressing any equipment
problem (spurious emissions, excessive splatter, an unstable signal)
brought to their attention by other operators.

\begin{quickreview}
\item Emission designators (e.g., J3E, F3E, A1A) internationally
describe bandwidth, modulation, and information type in a standardized,
country-independent code.
\item Simplex (shared frequency, taking turns), half-duplex (both
directions possible, not simultaneous from the operator's side), and
full duplex (simultaneous both ways) describe communication timing
structure.
\item Repeaters extend range by receiving on one frequency and
retransmitting on another (the offset), often gated by CTCSS/DCS tones.
\item Standard procedural conventions (CQ, roger/over/out, break-in)
exist to keep communication unambiguous across languages and
conditions.
\item Minimizing interference through good operating habits is a
universal amateur responsibility, independent of any specific national
rule.
\end{quickreview}

\begin{practiceproblems}
\item Identify the modulation type, and information type, encoded in
the emission designator \texttt{F3E}.
\item Explain the difference between half-duplex and full-duplex
operation, using a repeater as an example.
\item Explain the purpose of CTCSS tone squelch on a repeater.
\item Explain why listening before transmitting (particularly before
calling CQ) is considered good operating practice.
\item Explain why standardized procedural words like ``roger'' and
``over'' are preferred over improvised phrasing during radio contacts.
\end{practiceproblems}
